\begin{ex}
Mỗi ngày mới bắt đầu từ mấy giờ đêm (rạng sáng)?
\shortans{0}
\loigiai{
Một ngày mới bắt đầu từ 0 giờ đêm (12 giờ đêm).
\textbf{Đáp số:} $0$.
}
\end{ex}

\begin{ex}
Bố đi làm bằng xe máy lúc 7 giờ sáng và đến công ty lúc 8 giờ sáng. Bố đi trên đường mất bao nhiêu giờ?
\shortans{1}
\loigiai{
Thời gian đi đường là: $8 - 7 = 1$ (giờ).
\textbf{Đáp số:} $1$.
}
\end{ex}

\begin{ex}
Buổi trưa, cả nhà em bắt đầu ăn cơm lúc 11 giờ và ăn xong lúc 12 giờ trưa. Bữa cơm trưa kéo dài bao nhiêu giờ?
\shortans{1}
\loigiai{
Bữa cơm kéo dài: $12 - 11 = 1$ (giờ).
\textbf{Đáp số:} $1$.
}
\end{ex}

\begin{ex}
Giờ ngủ trưa ở trường bán trú bắt đầu từ 12 giờ trưa đến 2 giờ chiều. Các bạn được ngủ trưa trong bao nhiêu giờ?
\shortans{2}
\loigiai{
Thời gian ngủ trưa là: $2 - 0 = 2$ (giờ) (từ 12h đến 2h chiều là 2 tiếng).
\textbf{Đáp số:} $2$.
}
\end{ex}

\begin{ex}
Một tờ lịch tháng ghi: Ngày 5 là Thứ Hai. Hỏi ngày 12 cùng tháng đó là thứ mấy? (Thứ Hai: 2; Thứ Ba: 3).
\shortans{2}
\loigiai{
Vì $5 + 7 = 12$ nên ngày 12 cũng là Thứ Hai (ghi số 2).
\textbf{Đáp số:} $2$.
}
\end{ex}

\begin{ex}
Tờ lịch tháng ghi: Ngày 8 là Thứ Sáu. Hỏi ngày 15 cùng tháng đó là thứ mấy? (Thứ Sáu: 6; Thứ Bảy: 7).
\shortans{6}
\loigiai{
Vì $8 + 7 = 15$ nên ngày 15 cũng là Thứ Sáu (ghi số 6).
\textbf{Đáp số:} $6$.
}
\end{ex}

\begin{ex}
Nếu ngày 3 là Thứ Bảy thì ngày 10 cùng tháng đó là thứ mấy? (Thứ Bảy: 7; Chủ Nhật: 8).
\shortans{7}
\loigiai{
Vì $3 + 7 = 10$ nên ngày 10 cũng là Thứ Bảy (ghi số 7).
\textbf{Đáp số:} $7$.
}
\end{ex}

\begin{ex}
Nếu ngày 17 là Thứ Tư thì ngày 10 cùng tháng đó là thứ mấy? (Thứ Ba: 3; Thứ Tư: 4).
\shortans{4}
\loigiai{
Vì $17 - 7 = 10$ nên ngày 10 cũng là Thứ Tư (ghi số 4).
\textbf{Đáp số:} $4$.
}
\end{ex}

\begin{ex}
Đồng hồ reo lúc 6 giờ sáng. Nam thức dậy, vệ sinh cá nhân và ăn sáng mất 1 giờ. Hỏi Nam sẵn sàng đi học lúc mấy giờ?
\shortans{7}
\loigiai{
Nam sẵn sàng đi học lúc: $6 + 1 = 7$ (giờ sáng).
\textbf{Đáp số:} $7$.
}
\end{ex}

\begin{ex}
Chuyến tàu hỏa khởi hành từ ga lúc 8 giờ tối và chạy suốt đêm đến 6 giờ sáng hôm sau. Hỏi chuyến tàu chạy hết bao nhiêu giờ?
\shortans{10}
\loigiai{
Từ 8 giờ tối (20h) đến 12h đêm là 4 giờ. Từ 12h đêm đến 6h sáng là 6 giờ. Tổng cộng: $4 + 6 = 10$ (giờ).
\textbf{Đáp số:} $10$.
}
\end{ex}

\begin{ex}
Bạn Hùng có lịch tập bơi vào Thứ Ba, Thứ Năm và Thứ Bảy hàng tuần. Hỏi trong một tuần Hùng đi tập bơi bao nhiêu buổi?
\shortans{3}
\loigiai{
Có 3 buổi tập bơi: Thứ Ba, Thứ Năm, Thứ Bảy (tất cả 3 buổi).
\textbf{Đáp số:} $3$.
}
\end{ex}

\begin{ex}
Tháng có 4 tuần lễ trọn vẹn và thừa 2 ngày. Hỏi tháng đó có bao nhiêu ngày?
\shortans{30}
\loigiai{
4 tuần lễ là: $7 \times 4 = 28$ ngày. Thêm 2 ngày: $28 + 2 = 30$ (ngày).
\textbf{Đáp số:} $30$.
}
\end{ex}

\begin{ex}
Nếu hôm nay là ngày 14 thì 7 ngày trước là ngày mùng mấy?
\shortans{7}
\loigiai{
Bảy ngày trước là ngày: $14 - 7 = 7$.
\textbf{Đáp số:} $7$.
}
\end{ex}

\begin{ex}
Nếu hôm nay là ngày 18 thì 7 ngày sau là ngày bao nhiêu?
\shortans{25}
\loigiai{
Bảy ngày sau là ngày: $18 + 7 = 25$.
\textbf{Đáp số:} $25$.
}
\end{ex}

\begin{ex}
Lớp 1A đi tham quan dã ngoại lúc 7 giờ sáng và trở về trường lúc 4 giờ chiều (tức 16 giờ). Chuyến dã ngoại kéo dài bao nhiêu giờ?
\shortans{9}
\loigiai{
Từ 7 giờ sáng đến 4 giờ chiều là: $16 - 7 = 9$ (giờ).
\textbf{Đáp số:} $9$.
}
\end{ex}

\begin{ex}
Một ngày có 24 giờ. Em ngủ 9 giờ mỗi ngày. Hỏi thời gian em thức trong ngày là bao nhiêu giờ?
\shortans{15}
\loigiai{
Thời gian thức là: $24 - 9 = 15$ (giờ).
\textbf{Đáp số:} $15$.
}
\end{ex}

\begin{ex}
Nếu Thứ Hai là ngày 2 thì Chủ Nhật của tuần đó là ngày mùng mấy?
\shortans{8}
\loigiai{
Từ Thứ Hai đến Chủ Nhật có 6 khoảng ngày: $2 + 6 = 8$.
\textbf{Đáp số:} $8$.
}
\end{ex}

\begin{ex}
Nếu Chủ Nhật là ngày 9 thì Thứ Hai của tuần đó là ngày mùng mấy?
\shortans{3}
\loigiai{
Thứ Hai đầu tuần đó là ngày: $9 - 6 = 3$.
\textbf{Đáp số:} $3$.
}
\end{ex}

\begin{ex}
Một tiết học kéo dài khoảng 35 phút. Giờ ra chơi kéo dài khoảng 15 phút. Tổng thời gian một tiết học và một giờ ra chơi là bao nhiêu phút?
\shortans{50}
\loigiai{
Tổng thời gian là: $35 + 15 = 50$ (phút).
\textbf{Đáp số:} $50$.
}
\end{ex}

\begin{ex}
Đồng hồ điện tử hiện số: 08:00. Lúc đó là mấy giờ đúng?
\shortans{8}
\loigiai{
08:00 tương ứng là 8 giờ đúng.
\textbf{Đáp số:} $8$.
}
\end{ex}

\begin{ex}
Đồng hồ điện tử hiện số: 17:00. Lúc đó là mấy giờ chiều?
\shortans{5}
\loigiai{
17 giờ tương ứng với 5 giờ chiều (vì $17 - 12 = 5$).
\textbf{Đáp số:} $5$.
}
\end{ex}

\begin{ex}
Đồng hồ điện tử hiện số: 20:00. Lúc đó là mấy giờ tối?
\shortans{8}
\loigiai{
20 giờ tương ứng với 8 giờ tối (vì $20 - 12 = 8$).
\textbf{Đáp số:} $8$.
}
\end{ex}

\begin{ex}
Tháng Hai năm thường có 28 ngày. Tháng Hai năm thường có bao nhiêu tuần lễ trọn vẹn?
\shortans{4}
\loigiai{
Vì $4 \times 7 = 28$ nên có đúng $4$ tuần lễ trọn vẹn.
\textbf{Đáp số:} $4$.
}
\end{ex}

\begin{ex}
Bé được nghỉ tết trong 1 tuần lễ và 3 ngày. Hỏi bé được nghỉ tết tất cả bao nhiêu ngày?
\shortans{10}
\loigiai{
Tổng số ngày nghỉ tết là: $7 + 3 = 10$ (ngày).
\textbf{Đáp số:} $10$.
}
\end{ex}

\begin{ex}
Kim giờ quay được 1 vòng quanh mặt đồng hồ (từ số 12 đến số 12) hết bao nhiêu giờ?
\shortans{12}
\loigiai{
Kim giờ quay đúng 1 vòng mặt đồng hồ hết $12$ giờ.
\textbf{Đáp số:} $12$.
}
\end{ex}
